\documentclass[11pt]{article}

\usepackage[margin=1in]{geometry}
\usepackage{amsmath,amssymb}
\usepackage{graphicx}
\usepackage{subcaption}
\usepackage{bm}
\usepackage{url}
\usepackage{xcolor}
\usepackage{booktabs}
\usepackage{hyperref}

\graphicspath{{images/}{artifacts/figures/}}

% placeholders for values not yet measured
\newcommand{\res}[1]{{\color{blue}\textbf{[#1]}}}
\newcommand{\figph}[2][0.26\textheight]{%
  \fbox{\parbox[c][#1][c]{0.92\textwidth}{\centering\textsc{Placeholder figure.}\\[4pt]\small #2}}}

\title{%
Planning in the Real Game Tree:\\
A Continuous-Action EfficientZero-Style Graph Transformer for Mixed Doubles Curling
}

\author{}
\date{}

\begin{document}
\maketitle

\begin{abstract}
Curling is a continuous-action, execution-noisy, zero-sum sport whose mixed
doubles ends open from two pre-placed stones and resolve after ten alternating
throws. We learn a single closed-loop policy for it by combining ideas from
MuZero \cite{schrittwieser2020mastering} and EfficientZero
\cite{ye2021mastering} with the practical reality that curling has a fast,
calibrated physics simulator. A graph transformer shared across an
EfficientZero-style head set, namely a full-covariance mixture policy over the
release parameters $[v,\theta,\omega,y_0]$, a Gaussian value head over the
end-score differential, and an action-conditioned latent-dynamics head, is
trained with a BYOL-style consistency objective \cite{grill2020bootstrap}.
The learned dynamics is used only to shape the representation; all search and
all training targets use the real simulator. Continuous-action search produces
value-weighted policy targets, candidate values are averaged over a fitted
Student-$t$ execution-noise model \cite{yee2016monte} so the distilled policy
prefers shots that survive imperfect execution, and human priors are recovered
from position-only match data with a cross-entropy method inverse solver
\cite{deboer2005tutorial}. To avoid the curriculum forgetting we observed in a
naive horizon-by-horizon schedule, we collect per-horizon search targets once
(including mode-balanced canonical opening states that are absent from human
data) and \emph{consolidate} a single model on the union in one joint
training pass. Under realistic execution noise in true alternating play, the
consolidated model beats the human prior across every horizon, with the
largest gains at long horizons where strategic structure matters most.
\end{abstract}

\section{Introduction}

Curling is a precision sport with continuous actions, execution noise, and
multi-horizon tactics. Mixed doubles adds a fixed structure: each end opens
from two pre-placed stones and runs for ten alternating throws, so an end is a
ten-step episode whose outcome turns on a handful of decisive decisions. We
denote by \emph{horizon} $h \in \{1, \ldots, 10\}$ the number of throws
remaining at a given decision, including the current one. The opening shot
has $h=10$ and the final stone has $h=1$. Since throws strictly alternate
between the two teams, horizon parity also fixes the \emph{hammer}: the team
with the last stone (the hammer team) is decided by whether $h$ is odd or
even relative to which team throws first. We seek a closed-loop policy that
maps a board state and game context to an intended throw and is robust to
execution noise.

Modern model-based reinforcement learning, exemplified by MuZero
\cite{schrittwieser2020mastering} and EfficientZero \cite{ye2021mastering},
learns a latent dynamics model and plans inside it. That design is powerful
when the true dynamics are unknown or expensive, but it couples planning
quality to model fidelity, and learned models drift over long rollouts,
precisely the regime curling demands. We make a different choice. Curling has
a fast, calibrated, differentiable physics simulator, so treating it as the
transition removes model bias from search entirely. We therefore keep the
machinery that makes EfficientZero representations strong, namely a shared
encoder, an action-conditioned latent-dynamics head, and a BYOL-style
latent-consistency objective \cite{grill2020bootstrap}, but use it only to
shape the representation, never to roll out plans. All search and all
training targets use the real simulator.

Our contributions are:
\begin{enumerate}
\item \textbf{A simulator-grounded world model.} A single graph-transformer
trunk feeds a small, gated head set (policy, value, latent dynamics, and an
auxiliary two-step-return head). The simulator is the sole transition for
planning and targets, and the latent-dynamics head is a pure representation
auxiliary trained against a BYOL-style target encoder. We do not use a
per-step outcome distribution or a behavior-cloning loss in the iterative
phase.

\item \textbf{Continuous-action search with execution-noise robustness.} A
multi-ply KR-UCT tree \cite{yee2016monte} against the simulator, and a
cheaper one-ply robust selector that averages each candidate's decision
value over $K$ fitted-Student-$t$ noisy realizations. Both yield value-weighted
policy targets, and we ablate them head-to-head.

\item \textbf{Inverse-recovered human priors via the cross-entropy method.}
Release parameters are unobserved in match data, so we recover them from
observed stone displacements with a CEM \cite{deboer2005tutorial} inverse
solver and use the quality-filtered estimates to warm-start the trunk and
policy.

\item \textbf{Anti-forgetting consolidation.} Naive sequential per-horizon
training forgets late-horizon skill once the curriculum advances. We instead
collect per-horizon search targets once (including mode-balanced canonical
pre-placed states for the opening shot, which are absent from human data)
and consolidate a single model on the union of all per-horizon buffers in one
joint training pass.

\item \textbf{A faithful evaluation protocol.} We measure game strength by
true alternating play under both throwing orders, with execution noise on
realized throws, and report per-horizon \emph{with-hammer} win rate and score
differential, plus odd/even pair averages that cancel the hammer to isolate
skill.
\end{enumerate}

\section{Related Work}

\paragraph{Strategy and decision-making in curling.}
A growing body of work studies strategic decision-making in curling using
machine learning and search. Xiao et al.\ propose a policy decision system
designed for real competition scenes, combining a situation-aware perception
network with a digital extraction module and a domain-specific Curling MCTS
that plans in a continuous action space \cite{xiao2023policy}. Brunner and
Riesen compare supervised learning methods for position analysis in
simulated curling and show that neural and graph-based models outperform
heuristic baselines, while emphasizing that existing datasets remain
insufficient to consistently surpass professional curlers, particularly in
complex mid-end situations \cite{brunner2024comparing}. Han et al.\ combine
Neural Fictitious Self-Play with KR-UCT to approximate Nash strategies in
digital curling, focusing on large game-tree search in continuous action
spaces and automatic generation of supervised data \cite{han2022game}. Lee
et al.\ present a self-play deep reinforcement learning framework for
simulated curling that integrates continuous action search via kernel
regression with high-fidelity simulation, achieving strong performance in
an international digital curling competition without hand-crafted features
\cite{lee2018deep}. Our work shares the goal of high-quality continuous-action
shot selection and, like Han et al.\ and Lee et al., uses KR-UCT, which was
introduced for continuous-action MCTS with execution uncertainty by Yee
et al.~\cite{yee2016monte}. It differs in pairing search with an
EfficientZero-style graph-transformer whose shared representation is trained
by BYOL-style latent consistency, in planning directly against a calibrated
physics simulator rather than a learned model, in optimizing robustness to a
fitted Student-$t$ noise model, and in recovering human priors by inverse
dynamics from position-only match data.

\paragraph{Robotic and real-world curling.}
Beyond purely digital agents, several systems close the loop with the
physical world. Won et al.\ introduce Curly, a curling robot powered by an
adaptive deep reinforcement learning framework that incorporates temporal
features to handle nonstationary ice conditions
\cite{doi:10.1126/scirobotics.abb9764}. Curly attains human-level performance
in official matches against elite human teams, demonstrating that carefully
adapted RL can narrow the sim-to-real gap in curling. Xiao et al.\ also
emphasize deployment in real training environments, using perception modules
and tailored MCTS to bridge from image observations to strategic
recommendations \cite{xiao2023policy}. Compared to these systems, our work
remains fully in simulation but focuses on the complementary problem of
how to combine a learned graph-transformer representation with a calibrated
simulator for tree search over full ends.

\paragraph{Local tactical optimization in curling.}
Several papers examine specific subproblems rather than full-game strategy.
Ahmad et al.\ focus on the hammer shot, the final shot of an end, and
investigate algorithms for selecting an optimal action in a continuous,
low-dimensional, stochastic setting \cite{ahmad2016action}. They adapt a
Delaunay triangulation method to this setting and show that it can
outperform alternative algorithms and, under some conditions, even
Olympic-level human performance in their simulator. Our approach instead
targets multi-shot planning over an entire end, with continuous-action
KR-UCT search against a physics simulator and a learned value, and, like
Ahmad et al., optimizes directly in the continuous release-parameter space
$[v,\theta,\omega,y_0]$ rather than over a discrete library of high-level
shot types.

\paragraph{Physics-based strategy in other sports.}
The idea of using physics models or simulators to inform sports strategy
is not unique to curling. Fragkiadaki et al.\ learn visual predictive models
of physics in a simulated billiards domain and use them to plan goal-directed
actions via ``visual imagination'' over object-centric latent states
\cite{fragkiadaki2015learning}. In ten-pin bowling, Ji et al.\ derive a
rigid-body model for ball motion and use simulation to map initial conditions
to strike probability, revealing that imperfect bowlers may benefit from
particular targeting strategies that exploit the oil pattern to create
additional ``miss room'' \cite{ji2022using}. Conceptually we share these
works' use of an explicit physics model for planning, but rather than
\emph{learning} the dynamics for use inside the planner, we treat a
calibrated curling simulator as the transition and learn only a
representation that imitates it. We operate on object-centric stone positions
rather than pixels and embed the simulator inside a continuous-action search
loop for a multi-agent, zero-sum game.

\paragraph{World models and model-based reinforcement learning.}
Our method is also related to the broader literature on world models in
reinforcement learning. DreamerV3 \cite{hafner2023mastering} demonstrates that
a single world-model-based algorithm can, with one configuration, achieve
state-of-the-art performance across over 150 diverse tasks, including complex
3D environments, by learning a latent dynamics model and improving policies
through imagination. Alonso et al.\ introduce DIAMOND, a diffusion-based
world model for Atari that shows how modeling high-fidelity visual details
via diffusion can translate into improved downstream policy performance
\cite{alonso2024diffusion}. In our setting we draw on these ideas for
representation learning, namely an action-conditioned latent-dynamics head
and a BYOL-style latent-consistency objective \cite{grill2020bootstrap} in
the EfficientZero spirit \cite{ye2021mastering}, but, unlike Dreamer and
DIAMOND, we do \emph{not} roll out the learned model for planning. We operate
in a low-dimensional, object-centric state space and use a calibrated physics
simulator as the transition, reserving the learned dynamics for shaping the
shared graph-transformer representation.

\paragraph{Planning with learned models.}
Finally, our planning setup is inspired by MuZero
\cite{schrittwieser2020mastering}, which combines tree search with a learned
model to achieve superhuman performance in Atari, Go, chess, and shogi
without explicit knowledge of the game rules. MuZero learns a latent
dynamics model that produces policy, value, and reward predictions suitable
for search. In contrast, we deliberately do \emph{not} plan in the learned
latent model. MuZero learns dynamics, reward, value, and policy expressly
for in-model search, whereas we borrow only its shared-representation and
latent-consistency machinery and run the actual search against a calibrated
simulator with a continuous-action KR-UCT tree \cite{yee2016monte} and an
inverse-recovered human prior. Relative to prior curling work, which either
relies on a hand-coded simulator with classical search or on model-free RL
in continuous action spaces, our contribution is to combine an
EfficientZero-style graph-transformer representation with simulator-grounded
planning and execution-noise-robust, search-distilled policies for mixed
doubles curling.

\section{Problem Setting}
\label{sec:problem}

\paragraph{State, condition, action.}
We represent a board by the \textbf{state} $s\in[0,1]^{24}$: twelve stone
slots (six per team), each an $(x,y)$ position normalised by
$\textsc{pos\_max}=4095$ raw units ($1\,\text{raw}=0.003048$\,m). Empty or
out-of-play slots take a sentinel near $1$, and a slot is \emph{in play} iff
its raw coordinates lie strictly inside the sheet bounds. Game context is a
3-vector \textbf{condition}
$c=[\,\text{shot\_norm},\,\text{team\_order},\,\text{stone\_block}\,]$, where
$\text{shot\_norm}\in[0,1]$ encodes progress through the end,
$\text{team\_order}\in\{0,1\}$ indicates the throwing team's order (which
flips with each thrown stone), and $\text{stone\_block}\in\{0,1\}$ identifies
which six slots the thrower owns. The \textbf{action} is the continuous
4-vector $a=[\,v,\theta,\omega,y_0\,]$: release speed $v$ (m/s), aim angle
$\theta$ (rad), release spin $\omega$ (rad/s), and lateral release offset
$y_0$ (m). Action bounds are physically realistic and \emph{limit the search
and sample space} during planning (Appendix~\ref{app:sim}).

\paragraph{Horizon, hammer, episode.}
The \textbf{horizon} $h\in\{1,\ldots,10\}$ is the number of throws remaining
at a decision, including the current one. The opening shot has $h=10$ and
the final stone has $h=1$. The first thrown stone has $h=10$, the second
$h=9$, and so on through $h=1$. Because throws strictly alternate, horizon
parity determines the \textbf{hammer} (the team that throws the last stone):
$h$ odd at a decision means the team-to-move throws an odd-indexed remaining
stone, which puts the hammer on the team-to-move at odd horizons and on the
opponent at even horizons. An \textbf{end} is the full 10-throw episode
opened from the two pre-placed stones, and a curling \emph{game} (out of
scope here) consists of several ends; we evaluate on single ends.

\paragraph{Transition and scoring.}
A throw is realized by the curling physics simulator (Section~\ref{sec:sim};
Appendix~\ref{app:sim}), which integrates a friction and curl model with
smooth penalty-based stone-stone contacts and assigns the new stone to the
thrower's next free slot. Execution noise is modelled by perturbing the
intended action, $\tilde a = a+\epsilon$, with $\epsilon$ drawn from a fitted
heavy-tailed Student-$t$ (Section~\ref{sec:noise};
Appendix~\ref{app:noise}). After all ten throws the \textbf{end score} (the
\emph{ValueDiff}) is the signed margin $R(s)$ under standard curling rules
(only the team with the stone closest to the button scores, counting all of
its stones nearer than the opponent's closest), expressed in a chosen team's
perspective.

\paragraph{Free-guard-zone legality.}
Early-end takeouts are restricted: an opponent's rock may not be removed
from play until the fourth thrown stone of the end (active for thrown stones
1, 2, 3, i.e.\ horizon $h\ge 8$). Contact that keeps an opponent rock in
play is legal, as is removing one's own rock; only knocking an opponent
rock out of play is illegal, and the protection covers every opponent rock
in play, including the pre-placed one. When a throw violates the rule, the
delivery is forfeited and the board reverts to the pre-shot configuration.
We apply this legality operator to every simulated post-state used in search
and target generation (Appendix~\ref{app:fgz}).

\paragraph{Objective.}
We seek a closed-loop policy $\pi(a\mid s,c)$ that maximises the expected
terminal end margin under execution noise, played as a two-player zero-sum
game in which the opponent uses the same (current) policy.

\section{Method}
\label{sec:method}

\subsection{Overview}
A single graph-transformer encoder $E(s,c)\to h\in\mathbb{R}^{256}$ feeds
four gated heads (Figure~\ref{fig:arch}): a policy head $\pi(\cdot\mid h)$,
a value head $V(h)$, an action-conditioned latent-dynamics head $G(h,a)$,
and an auxiliary two-step-return head $\hat r_2(h)$ used in selection
variants. Following MuZero \cite{schrittwieser2020mastering} and
EfficientZero \cite{ye2021mastering}, we form a $K$-step latent unroll
$h_0=E(s_0,c_0)$, $h_{k+1}=G(h_k,a_k)$, and predict policy and value at
every latent step. Unlike MuZero and EfficientZero, the unrolled latent is
never used to produce transitions for planning or for value/policy targets:
the simulator is the transition. The dynamics head is supervised only by a
latent-consistency loss against a BYOL-style \cite{grill2020bootstrap} target
encoder of the simulator's true next state, so it learns to imitate the
encoded physics and thereby regularises the shared trunk
(Section~\ref{sec:targets}). We do not use a per-step outcome distribution
and we do not use a behavior-cloning loss in the iterative phase; the human
prior enters only through the warm start (Section~\ref{sec:inverse}).

\begin{figure}[t]
    \centering
    \figph[0.30\textheight]{Architecture overview. The shared GraphTF trunk
    $E(s,c)\to h$ (22 nodes: 12 stone-slot nodes, 9 geometric landmark nodes,
    and 1 condition-carrying global token; edge-biased attention with
    9-dimensional edge features) feeds the policy (full-covariance Gaussian
    mixture over $[v,\theta,\omega,y_0]$), value (Gaussian over end margin),
    latent dynamics $G(h,a)$, and auxiliary two-step-return $\hat r_2$ heads.
    Dashed path: BYOL-style latent consistency
    $G(E(s),a)\approx \mathrm{sg}\,E_{\bar\theta}(\mathrm{Sim}(s,a))$ against
    a slow-moving target encoder, representation only, never used for
    planning.}
    \label{fig:arch}
\end{figure}

\subsection{Shared graph-transformer trunk}
\label{sec:trunk}
The encoder constructs an object-centric graph at every step so that the
model is permutation-equivariant over the twelve stone slots and can attend
selectively to relevant board geometry rather than processing a fixed
feature vector.

\paragraph{Nodes (22).} The graph has 22 nodes: 12 \emph{stone-slot} nodes,
9 \emph{landmark} nodes that expose drawing and takeout geometry, and 1
\emph{global} token that carries the condition $c$ and serves as the
trunk's readout. The 9 landmarks consist of the button centre, three
release-area reference points across the nearer hog line (left, centre,
right, at $\approx 6.4$\,m forward of the button), and five behind-house
takeout-boundary points placed $\approx 2.4$\,m behind the button (just
past the back of the house) at lateral offsets $\{-1.83, -0.91, 0, +0.91,
+1.83\}$\,m. Each node has a 6-dimensional feature vector with fields
$[x, y, \text{team\_indicator}, \text{is\_live}, \text{is\_landmark},
\text{spare}]$, where \texttt{spare} is unused under the current
configuration. The global token is prepended at encode time (so $h$ is
recovered as the first row of the encoder output) and its features carry
the condition projected to the hidden dimension.

\paragraph{Edges (9-d features).} Edges are dense (all node pairs). Each
edge carries a 9-dimensional feature vector: 4 base geometric scalars
$[\Delta x, \Delta y, \text{distance}, \text{same\_team\_indicator}]$ plus 5
curling-specific scalars:
\begin{enumerate}\setlength\itemsep{0pt}
\item[(1)] \textbf{scorability}, the unoccluded angular span of the target
stone as seen from the button (broadcast as a node-level feature into each
incident edge);
\item[(2)] \textbf{reach}, the feasibility of physically reaching the target
along a curl arc on release-source edges, or pairwise angular reachability
otherwise;
\item[(3)] \textbf{score-out compatibility}, how compatible the
outgoing-direction is with drawing to the button;
\item[(4)] \textbf{takeout-out compatibility}, how compatible the
outgoing-direction is with ejecting an opponent through the behind-house
corridors;
\item[(5)] the multiplicative composite
$\textsc{reach}\cdot\max(\textsc{score-out},\textsc{takeout-out})$.
\end{enumerate}
These scalars are recomputed per state from the live JAX simulator's
geometry rather than learned. Edges are not pruned: the attention layers
see all node pairs.

\paragraph{Edge-biased attention.} The trunk is a 4-layer transformer of
hidden width $256$ with $4$ attention heads and dropout $0.1$
($\approx 7.3\text{M}$ parameters total). Each layer applies an
\emph{edge-biased} multi-head attention in which the edge feature for the
pair $(i,j)$ is projected to (a) a per-head additive bias added to the
attention logits, and (b) a value modulation that gates the attended values
from $j$ to $i$. The block then applies a residual feed-forward (4$\times$
expansion) with layer normalisation. The latent $h$ is read out as the
output embedding of the global token. Inputs are augmented at training time
by a horizontal flip across the centerline combined with a team-slot swap
that exchanges stones $1$--$6$ and $7$--$12$ and flips
$\text{stone\_block}$; the flip negates the angle, spin, and lateral-offset
components of all action targets accordingly, so policy, value, dynamics,
and consistency targets remain consistent under both operations.

\paragraph{Warm start.} The trunk is initialised from the human-prior policy
network's matched modules (the global-token vector, node and condition
projections, and transformer layers), and the value head is initialised from
a separately trained standalone value model that uses the same graph
configuration. The BYOL target encoder is then synced to the trunk before
training begins. The policy head is optionally initialised from the prior
(true by default) and is then refined by search distillation.

\subsection{Heads}
\label{sec:heads}
\textbf{Policy.} $\pi(\cdot\mid h)$ is a full-covariance Gaussian mixture
($K=16$) over the standardised action $z=(a-\mu_a)/\sigma_a$, parameterised
by mixture logits, means, and Cholesky factors. We report the negative
log-likelihood $\mathrm{NLL}_\pi$ under this density.

\textbf{Value.} $V(h)=(\mu_V,\log\sigma_V^2)$ is a Gaussian over the signed
end margin (the ValueDiff). We train it on \emph{realized} terminal margins
only, not on search statistics (Section~\ref{sec:targets}).

\textbf{Latent dynamics.} $G(h,a_{\text{box}})$ maps the current latent and
a box-normalised action $a_{\text{box}}\in[-1,1]^4$ to a predicted next
latent via a residual MLP, used only for latent unrolling and the
consistency loss.


\subsection{Simulator}
\label{sec:sim}
The transition $s'=\mathrm{Sim}(s,c,a)$ is a JAX physics integrator using
semi-implicit Euler: macro step $dt=0.02$\,s, two substeps, up to
$1{,}500$ steps, terminating when all stone speeds fall below $0.03$\,m/s.
Dynamics combine linear ice friction (deceleration $\approx
0.11\,\text{m/s}^2$ opposing motion), a spin-induced curl lateral
acceleration with strength $k_{\text{curl}}\approx 0.10$ and a
speed-saturated factor capped at $c_{\text{cap}}=2.5\,\text{m/s}$,
exponential spin decay (rate $\approx 0.15\,\text{s}^{-1}$ with additional
damping during contact), and smooth penalty contact forces with stiffness
$\sim 2.5\times10^4$\,N/m, damping $\sim 220$\,kg/s. Stone radius is
$0.145$\,m and stone mass is $19.96$\,kg. Geometry is full-sheet
($d_{\text{hog-tee}}=28.35$\,m). Numerical parameters, action bounds, and
the contact model are listed in Appendix~\ref{app:sim}. Crucially, no calls
to the learned dynamics head enter search or target generation: the
simulator is the only transition we plan against. The JAX implementation is
batched and runs on CPU or GPU.

\subsection{Continuous-action search and target generation}
\label{sec:search}
We collect search-improved policy targets at a set of MCTS \emph{root}
states $\{s,c\}$. ``Root'' here means an MCTS root: a board position at which
search is run to produce an improved policy at that single state, which is
then stored as a training target. Roots are drawn from real human boards at
each throws-remaining horizon $h\in\{1,\ldots,9\}$, plus canonical pre-placed
openings at $h=10$ (see Section~\ref{sec:phases}). For each root, the search
produces a finite set of candidate actions $\{a_m\}_{m=1}^{M}$ together with
values $\{q_m\}$ expressed in the root team's perspective.

\paragraph{Multi-ply KR-UCT.}
Following Yee, Lis\'y, and Bowling \cite{yee2016monte}, a recursive tree
against the simulator combines two ingredients that are essential in
continuous-action MCTS. First, \emph{progressive widening}: a node gains a
new child only when its visit count is large enough
($|\text{children}|\le k\,N^{\alpha}$), with each new action sampled from
the current neural policy and de-duplicated in action space. Second,
\emph{kernel-regression UCT selection}: among existing children, each
child's mean value is smoothed by kernel regression over its neighbours in
$z$-normalised action space (Gaussian kernel of fixed bandwidth), with the
UCT exploration bonus folded into the kernel mass; the argmax of the
smoothed scores is selected. This shares statistics across nearby continuous
actions, which is essential when each action is visited few times. A
freshly-expanded leaf is evaluated by rolling forward to the end of the end
under the simulator, using the \emph{search-distilled (or, at iteration
zero, human-prior) policy itself as the rollout policy}, i.e.\ the rollout
is on-policy, not uniform random. The terminal board is then scored by the
curling rules.

\paragraph{Cheaper one-ply robust selector.}
The deep tree is expensive. As a cheap alternative we use a one-ply root
bandit in the spirit of \cite{yee2016monte}: sample $M$ candidate actions
from the policy and score each candidate by the \emph{mean} of its value
over $K$ noisy executions $\tilde a = a+\epsilon$ (or, in the $\hat r_2$
variant, by $-\hat r_2$ on each post-state). The averaging is the entire
robustness mechanism: selection prefers shots whose decision value
\emph{survives} execution error rather than brittle perfect-line shots. We
ablate this one-ply robust selector against the deeper KR-UCT tree
(Section~\ref{sec:exp}); the one-ply selector wins our deployable
comparison, and is what fed the consolidated model below.

\paragraph{Policy target as a value-weighted set of searched actions.}
Both selectors return $\{(a_m, q_m)\}$. We convert this to a target with a
\emph{value-weighted soft-top-$k$} over the searched actions: keep the
top-$M$ candidates by $q$, with weights $w_m \propto \exp(q_m/\tau)$ summing
to $1$. Intuitively, the policy is trained to place probability mass on the
high-value candidates the search found, weighted continuously by how good
each one looked, rather than by visit count as in discrete AlphaZero/MuZero.
The training signal is a weighted mixture NLL of these candidate actions
under $\pi(\cdot\mid h_0)$.

\paragraph{Value target as the realized terminal margin.}
The value head is trained on the \emph{realized} signed end margin obtained
by rolling the policy to the terminal state under the simulator and scoring
under the curling rules: an unbiased, terminal-grounded Monte-Carlo return.
We do not regress the value head onto search Q-values, which it overfits in
practice.

\subsection{Execution-noise model}
\label{sec:noise}
Execution noise on $a=[v,\theta,\omega,y_0]$ is a diagonal Student-$t$
($\nu=5$, heavier-tailed than Gaussian). The \emph{speed} and \emph{angle}
scales are the maximum-likelihood fits of Yee et al.~\cite{yee2016monte}
to the World Curling Federation's \emph{Draw Shot Challenge} dataset,
the canonical real-world source for curling execution error and the same
calibration used in their KR-UCT experiments. The angle scale is
speed-dependent (larger at lower speeds). For \emph{spin} and \emph{lateral
offset} no comparable public calibration data exists, so we use small
empirical constants. This is a known limitation: under-modelling spin and
offset noise could leave the policy over-confident in shots whose accuracy
depends sensitively on those parameters, and a calibrated spin/offset noise
model is flagged as future work. The same noise model is used both to
perturb realized transitions during search and evaluation and to average
candidate values during target generation, so the policy is optimised
end-to-end for robustness to it.

\subsection{Inverse-recovered human prior}
\label{sec:inverse}
The match data records stone positions before and after each shot and a
coarse task label, but not the continuous release parameters. We recover
them with an inverse-dynamics solver based on the Cross-Entropy Method
(CEM) \cite{deboer2005tutorial}, the standard gradient-free importance-
sampling optimiser for noisy objectives. At each iteration we sample a
population of candidate actions from a Gaussian, simulate each forward,
score against the observed post-state with a permutation-invariant
team-ownership-aware position loss that accounts for stones knocked out of
play, and refit the Gaussian to the elite fraction. We keep only confidently
solved, low-residual estimates (the ``realistic'' subset) and use them to
train the human prior, which warm-starts the trunk and policy/value heads.
The inverse-recovered actions are \emph{not} used as a behavior-cloning loss
during the iterative phase (Appendix~\ref{app:inverse}).

\subsection{Targets and losses}
\label{sec:targets}
For each collected root we run the search above and form a fixed-shape
record. The training loss is a masked sum over the $K$-step latent unroll:
\begin{equation}
\mathcal{L} = \lambda_{\pi}\mathcal{L}_{\text{distill}}
            + \lambda_{V}\mathcal{L}_{\text{value}}
            + \lambda_{C}\mathcal{L}_{\text{consist}}
            + \lambda_{r}\mathcal{L}_{\text{r2}} .
\label{eq:loss}
\end{equation}

\textbf{Policy distillation.} $\mathcal{L}_{\text{distill}}$ is the
weight-$w_m$ mixture NLL of the searched actions $\{a_m\}$ under
$\pi(\cdot\mid h_0)$ (Section~\ref{sec:search}).

\textbf{Value.} $\mathcal{L}_{\text{value}}$ is an MSE-plus-Gaussian-NLL
regression toward the realized terminal margin (Section~\ref{sec:search}).

\textbf{Latent consistency (BYOL-style).}
$\mathcal{L}_{\text{consist}}$ is a SimSiam-style loss between the unrolled
latent $h_{k}=G(h_{k-1},a_{k-1})$ and the stop-gradient encoding of the
simulator's true next state under a target encoder,
$\mathrm{sg}\,E_{\bar\theta}(\mathrm{Sim}(s_{k-1},a_{k-1}))$. The target
encoder $E_{\bar\theta}$ is an identical copy of the trunk whose weights
$\bar\theta$ are not trained by gradient descent but are tracked as a slow
\emph{moving average} of the online encoder's weights $\theta$, updated
after every optimisation step,
\begin{equation}
\bar\theta \;\leftarrow\; \tau\,\bar\theta + (1-\tau)\,\theta,
\qquad \tau = 0.99 .
\label{eq:ema}
\end{equation}
This is the BYOL device \cite{grill2020bootstrap}: averaging the weights
yields a slowly-moving, stable regression target rather than one that shifts
every step, which is what prevents the SimSiam-style consistency objective
from collapsing. This consistency term is the only signal that trains $G$,
and it is the mechanism by which world-model learning shapes the shared
representation \emph{without} the learned model ever being used to plan.


\subsection{Training: per-horizon collection, joint consolidation}
\label{sec:phases}
The naive way to train a multi-horizon model is a sequential curriculum:
collect at horizon $h$, train, advance to $h{+}1$, repeat. We tried this and
observed a clear forgetting pattern: extending into late horizons overwrites
earlier-horizon skill in the policy. Importantly, the value head does not
forget because it trains on a horizon-agnostic real terminal buffer; only
the per-horizon policy distillation drifts. We therefore split training
into two decoupled phases:

\paragraph{(1) Per-horizon target collection.}
For each horizon $h\in\{1,\ldots,10\}$ we collect a buffer of search-improved
records using the one-ply robust selector against the simulator, starting
from the human-prior warm start and refreshing the policy used inside the
search between stages so that later horizons benefit from earlier learning.
The roots at $h\le 9$ are real boards drawn from the human dataset filtered
by throws-remaining. \emph{The opening shot} ($h=10$) is not in the dataset
because annotators skipped the canonical mixed-doubles pre-placement; we
synthesise the canonical openings ourselves (\textsc{standard} and the two
power-play variants \textsc{pp\_left} and \textsc{pp\_right}, each with two
guard-slot assignments for a total of 6 distinct openings) and
\emph{stratify} the collection budget so each opening receives an equal
number of search roots. Without this balancing, proportional sampling from
the (skewed) data would leave the power-play openings under-collected; with
it, they are not.

\paragraph{(2) Consolidation.}
We then take the union of all per-horizon buffers and train a single model
on the mixture in one joint pass. Each minibatch is sampled uniformly across
horizons, so no horizon is overwritten. The consolidated model warm-starts
from the per-stage champion (the best per-horizon iterate), optimiser state
is reset (a fresh Adam moment estimate), and the value head continues to
train on the horizon-agnostic real terminal buffer. This is cheap: no
re-collection, one training-cost pass. Empirically it removes the forgetting
and exceeds the per-stage champion across every horizon
(Section~\ref{sec:exp}).

Collection (simulator-heavy, JAX) and training (PyTorch DDP) run as separate
processes to avoid GPU contention; details in Appendix~\ref{app:train}.

\section{Experiments}
\label{sec:exp}

\paragraph{Data.}
We use shot-by-shot records from four international mixed-doubles
competitions ($\approx 26{,}000$ board states), held out by competition for
value evaluation, plus synthetic terminal states for value calibration. The
human prior is trained on the inverse-recovered, quality-filtered action
subset (CEM solver-converged with refine residual $\le 0.08\,\text{m}^2$).

\paragraph{Metrics.}
We distinguish two kinds of metric.

\emph{Auxiliary metrics (vs.\ human data).} For the value head we report
mean squared error and Gaussian-NLL against realized end margins on the
held-out competition. For the policy we report human-action NLL. These
compare the trained heads to human behaviour and are not a direct measure
of game strength. A stronger policy can have higher policy-NLL precisely
because it concentrates on high-value, robust shots that diverge from the
broader human distribution. We report them only to track training health,
not to rank players.

\emph{Game strength (primary).} True alternating play to the terminal state
under realistic execution noise. Each decision is made by the player's own
one-ply robust selector ($M=48$ candidates, mean decision value over $K=8$
noisy executions), and the executed throw is a noisy realisation of the
intended action. We play both throwing orders per root and report
\textbf{per-horizon win rate with the hammer} and the matching
\textbf{without-hammer} numbers, the \textbf{score differential} for each,
and \textbf{odd/even pair averages}
$\big(\!\!\text{wr}(h)\!+\text{wr}(h{+}1)\big)/2$ that cancel the hammer
across two consecutive horizons to expose pure skill.

\paragraph{Four hammer-by-parity cases.}
Since throws strictly alternate, horizon parity fixes whether the team
\emph{to move now} also throws \emph{last}: at odd $h$ the team to move
holds the hammer; at even $h$ the opponent does. Crossing this with the
throwing-order swap our evaluation runs over (so that each player holds the
hammer in exactly one of the two orders per root) gives four meaningful
conditions for our model A:
\begin{enumerate}\setlength\itemsep{0pt}
\item[(a)] \textbf{with-hammer, odd $h$:} A is to move now \emph{and} throws
last. A is on the clock, and the final-stone advantage is on its side.
\item[(b)] \textbf{without-hammer, odd $h$:} the opponent is to move now and
will also throw last. A is waiting to throw and is at the hammer
disadvantage.
\item[(c)] \textbf{without-hammer, even $h$:} A is to move now but the
opponent throws last. A throws under pressure without the final-stone
safety net.
\item[(d)] \textbf{with-hammer, even $h$:} the opponent is to move now,
but A throws last. A is waiting to throw with the hammer in hand.
\end{enumerate}
Cases (a) and (c) are the ones in which A's own decision quality is
exercised at the root, since A is on the clock; (b) and (d) test how A
plays out a position where the opponent throws first. Column
\emph{with hammer} in Table~\ref{tab:main} reports (a) at odd $h$ and (d)
at even $h$; \emph{without hammer} reports (b) at odd $h$ and (c) at
even $h$. The pair averages
$\tfrac{1}{2}\big(\text{wr}(h)+\text{wr}(h{+}1)\big)$ are taken over the
A-as-to-move side at each horizon, i.e.\ they average case (a) at odd $h$
with case (c) at the adjacent even $h$. Both are ``A on the clock'', but
one of them gives A the hammer and the other does not, so the pair average
isolates skill from the hammer advantage.

\paragraph{Baselines.}
(i) the inverse-recovered human prior, (ii) the per-stage champion (best
single-horizon iterate before consolidation), (iii) a sequential
$h=1$-to-$h=10$ curriculum without consolidation (forgetting baseline),
and (iv) a one-ply deterministic decision-value selector (no noise
averaging).

\paragraph{Auxiliary metrics (val/test).}
The human prior model achieves val policy NLL $4.82$ on its own
inverse-recovered action validation set (selected at epoch 50). The
standalone human-prior value model achieves val MSE $2.13$ on the held-out
competition. Our consolidated full-end model has val policy NLL $19.97$ and
val value MSE $2.23$ on the same held-out splits. The policy-NLL gap is
expected and intended: our policy is not trained to imitate, and concentrates
mass on high-value robust shots. The value-MSE gap is small and not
predictive of game strength; see below.



\begin{table}[ht]
\centering
\caption{Aggregated pair averages and overall mean playing skill against the
human prior. Combining consecutive odd/even game lengths cancels the
situational advantage of the hammer to isolate pure strategic skill. Per-row
best values bolded across all three columns. NOISY alternating play, both
throwing orders, free-guard-zone rule, full val pool per horizon (no random
subsampling); see Table~\ref{tab:horizons} for per-horizon sample sizes.}
\label{tab:averages}
\small
\setlength{\tabcolsep}{8pt}
\renewcommand{\arraystretch}{1.1}
\begin{tabular}{r cc | cc | cc}
\hline
 & \multicolumn{2}{c|}{\textbf{Our Unified Model}} & \multicolumn{2}{c|}{\textbf{Per-Stage Baseline}} & \multicolumn{2}{c}{\textbf{Sequential Baseline}} \\
 \cline{2-7}
Setting & win & $\Delta$score & win & $\Delta$score & win & $\Delta$score \\
\hline
pair $(1,2)$  & $0.517$           & $+0.06$           & $\mathbf{0.525}$  & $\mathbf{+0.11}$  & $0.513$           & $+0.07$           \\
pair $(3,4)$  & $\mathbf{0.577}$  & $+0.29$           & $\mathbf{0.577}$  & $\mathbf{+0.36}$  & $0.561$           & $+0.10$           \\
pair $(5,6)$  & $\mathbf{0.552}$  & $\mathbf{+0.20}$  & $0.540$           & $+0.13$           & $0.515$           & $-0.08$           \\
pair $(7,8)$  & $\mathbf{0.537}$  & $-0.01$           & $0.521$           & $\mathbf{+0.05}$  & $0.527$           & $-0.12$           \\
pair $(9,10)$ & $\mathbf{0.576}$  & $\mathbf{+0.29}$  & $0.536$           & $+0.01$           & $0.491$           & $-0.19$           \\
\hline
\textbf{Mean Skill} & $\mathbf{0.552}$  & $\mathbf{+0.17}$  & $0.540$  & $+0.13$  & $0.521$  & $-0.04$ \\
\hline
\end{tabular}
\end{table}



\paragraph{Main result.}
Tables~\ref{tab:averages} and \ref{tab:horizons} report head-to-head game
strength vs.\ the human prior for the consolidated model (per-horizon
collection with mode-balanced canonical $h{=}10$ openings, then one joint
consolidation pass), the per-stage baseline (the best single-horizon
iterate, trained without consolidation), and a sequential
$h{=}1$-to-$h{=}10$ curriculum without consolidation. With the hammer the
consolidated model wins $0.70$ to $0.88$ of ends at every horizon with
strictly positive score differential ($+0.87$ to $+1.91$ points/end).
Without the hammer it loses by similar magnitudes in the opposite
direction; this is the expected hammer asymmetry and is what the
pair-average column of Table~\ref{tab:averages} cancels. On the
hammer-neutral pair average the consolidated model attains
\textbf{0.552} mean playing skill against the prior with \textbf{$+0.17$
points-per-end} score differential, beating the per-stage champion's
$0.540$ ($+0.13$) and the sequential curriculum's $0.521$ ($-0.04$). The
consolidated model wins or ties on win rate at every pair, and is the
only one of the three with a strictly positive score differential
averaged across the end. The sequential curriculum exhibits clear
forgetting: at pair $(9, 10)$ it drops to $0.491$ (below the prior) with
$\Delta=-0.19$, and its mean score differential is negative ($-0.04$
points/end). Per-end margins of this size compound across the eight ends
of a real game to a roughly one-point expected lead, which is a decisive
gap in mixed doubles practice.

\paragraph{Ablations.}
The head gating, search variants, and training phases let us isolate each
component on an otherwise identical pipeline:
\begin{itemize}
\item[(a)] \textbf{policy+value only} (no world model);
\item[(b)] \textbf{+latent dynamics with BYOL-style consistency}
(representation auxiliary on);
\item[(c)] \textbf{noise-robust vs.\ noise-free targets} (does noise
averaging during target generation actually help?);
\item[(d)] \textbf{sequential curriculum vs.\ consolidation} (does joint
training on the union of per-horizon buffers remove the forgetting we
observe in pure sequential training?);
\item[(e)] \textbf{KR-UCT search vs.\ vanilla UCT search} (does the kernel
regression over neighbouring continuous actions matter, or does naive UCB1
on raw candidates suffice?);
\item[(f)] \textbf{multi-ply KR-UCT targets vs.\ one-ply robust targets}
(is the deep tree worth the cost over the cheaper one-ply selector?);
\item[(g)] \textbf{proportional vs.\ mode-balanced pre-placed $h{=}10$
collection} (the openings are skewed in the data but synthetic in
collection, so we test whether a balanced collection budget helps
power-play openings).
\end{itemize}
\res{Ablation table.} The scientific question behind (b) is whether the
EfficientZero-style consistency objective improves the policy/value heads
through representation alone, given that the learned dynamics is never used
to plan. (e) and (f) settle the search-design questions.

\paragraph{Qualitative analysis.}
We visualise (i) full ends played by the consolidated policy against itself
and against the human prior, with realised throw trajectories drawn by the
physics simulator (Figure~\ref{fig:gameplay}), and (ii) policy proposal
heatmaps for representative board states: where the policy's sampled plans
would land on empty ice, isolating its intended targets from collisions
(Figure~\ref{fig:heatmaps}).

\begin{figure}[t]
    \centering
    \figph{Full-end games (pre-placement + ten throws): consolidated policy
    vs.\ itself (clean execution) and vs.\ the human prior under the noisy
    winrate setup. Trajectories are the actual simulated paths (curl plus
    collisions), not straight-line approximations.}
    \label{fig:gameplay}
\end{figure}

\begin{figure}[t]
    \centering
    \figph{Policy proposal heatmaps: density of where $\sim 3000$ sampled
    plans would land on empty ice for representative board states (the
    canonical pre-placed openings, plus mid-end positions at $h=6, 4, 2$).
    The actual board stones are overlaid; landing density isolates the
    policy's \emph{intended} targets from the collisions a real throw would
    produce.}
    \label{fig:heatmaps}
\end{figure}

\section{Limitations}
The transition is only as faithful as the calibrated simulator. Sim-to-real
gaps (ice variation, sweeping, player skill, opponent style) are not
modelled, and the inverse solver recovers a single most-likely action per
observed transition rather than a posterior, so human ``intent'' is
approximated. KR-UCT with on-policy Monte-Carlo leaves is accurate but
expensive, bounding the search budget per decision. The execution-noise
model is calibrated for \emph{speed} and \emph{angle} only (Draw Shot
Challenge data); \emph{spin} and \emph{lateral offset} noise are set to
small empirical constants in the absence of a public dataset, and a
calibrated spin/offset noise model is left as future work. Finally, the
learned dynamics head is deliberately not used for planning; whether a
sufficiently accurate learned model could safely prefilter candidates
(cheaply ranking proposals before simulation) is left to future work,
though the architecture exposes the hook for it.

\section{Conclusion}
We presented an EfficientZero-style graph-transformer for mixed doubles
curling that keeps the representation-learning machinery of model-based RL
(a shared encoder, an action-conditioned latent-dynamics head, and
BYOL-style latent consistency \cite{grill2020bootstrap,ye2021mastering})
while deliberately planning in the real game tree with a calibrated physics
simulator rather than in a learned latent model
\cite{schrittwieser2020mastering}. Continuous-action search (KR-UCT
\cite{yee2016monte} or the cheaper one-ply robust selector that wins our
ablation) supplies value-weighted policy-distillation targets, value targets
are realized terminal end scores, the human prior is recovered by a CEM
\cite{deboer2005tutorial} inverse solver, and execution-noise averaging
\cite{yee2016monte} yields robust shot selection under the correct
free-guard-zone rule. To avoid the forgetting we observed in sequential
per-horizon training, we collect per-horizon search targets once (with the
missing opening-shot states synthesised and mode-balanced) and consolidate
a single model on the union of all per-horizon buffers. The result is a
single full-end model that beats the human prior across every horizon under
realistic execution noise.

\appendix

\begin{table}[ht]
\centering
\caption{Horizon-by-horizon game performance of the three models against
the human prior under noisy alternating play with the corrected
free-guard-zone rule. Each horizon uses the full val pool with no random
subsampling: $h{=}1$--$9$ are $n_{\text{per\_order}}=95$--$133$ distinct
recorded human boards from the held-out competition, $h{=}10$ is $n=387$
canonical pre-placed openings frequency-weighted to the data distribution
across 9 distinct $(\text{mode},\,\text{guard\_slot},\,\text{thrower\_block})$
canonical states; both throwing orders played per root. \textbf{With Hammer}
indicates the throwing order in which our model holds the last-stone
advantage; \textbf{Without Hammer} indicates the opposite order. Per-row
best across the three models for each (winrate, $\Delta$score)$\times$(with,
without) cell is bolded.}
\label{tab:horizons}
\footnotesize
\setlength{\tabcolsep}{3pt}
\renewcommand{\arraystretch}{1.1}
\begin{tabular}{r cccc | cccc | cccc}
\hline
 & \multicolumn{4}{c|}{\textbf{Our Unified Model}} & \multicolumn{4}{c|}{\textbf{Per-Stage Baseline}} & \multicolumn{4}{c}{\textbf{Sequential Baseline}} \\
 \cline{2-13}
 & \multicolumn{2}{c}{With Hammer} & \multicolumn{2}{c|}{Without Hammer} & \multicolumn{2}{c}{With Hammer} & \multicolumn{2}{c|}{Without Hammer} & \multicolumn{2}{c}{With Hammer} & \multicolumn{2}{c}{Without Hammer} \\
 \cline{2-13}
$h$ & win & $\Delta$ & win & $\Delta$ & win & $\Delta$ & win & $\Delta$ & win & $\Delta$ & win & $\Delta$ \\
\hline
 1  & $0.704$           & $+0.87$           & $\mathbf{0.328}$  & $-0.76$           & $\mathbf{0.712}$  & $\mathbf{+0.90}$  & $0.304$           & $-0.81$           & $0.696$           & $+0.84$           & $0.324$           & $\mathbf{-0.73}$  \\
 2  & $0.722$           & $+1.00$           & $0.331$           & $-0.74$           & $\mathbf{0.774}$  & $\mathbf{+1.11}$  & $\mathbf{0.338}$  & $\mathbf{-0.68}$  & $0.714$           & $+1.05$           & $0.331$           & $-0.69$           \\
 3  & $\mathbf{0.814}$  & $+1.26$           & $0.297$           & $-0.93$           & $0.797$           & $\mathbf{+1.36}$  & $0.263$           & $-1.02$           & $0.801$           & $+1.21$           & $\mathbf{0.314}$  & $\mathbf{-0.78}$  \\
 4  & $\mathbf{0.706}$  & $+1.05$           & $0.339$           & $-0.69$           & $0.674$           & $\mathbf{+1.12}$  & $\mathbf{0.358}$  & $\mathbf{-0.65}$  & $0.670$           & $+0.79$           & $0.321$           & $-1.02$           \\
 5  & $\mathbf{0.839}$  & $\mathbf{+1.55}$  & $0.288$           & $-0.87$           & $0.780$           & $+1.30$           & $0.246$           & $-1.18$           & $0.737$           & $+0.92$           & $\mathbf{0.305}$  & $\mathbf{-0.79}$  \\
 6  & $0.832$           & $+1.62$           & $0.265$           & $-1.15$           & $\mathbf{0.867}$  & $\mathbf{+1.78}$  & $\mathbf{0.301}$  & $\mathbf{-1.05}$  & $0.770$           & $+1.13$           & $0.292$           & $-1.07$           \\
 7  & $0.828$           & $+1.38$           & $0.222$           & $-1.37$           & $\mathbf{0.848}$  & $\mathbf{+1.49}$  & $0.273$           & $-0.98$           & $0.798$           & $+1.10$           & $\mathbf{0.283}$  & $\mathbf{-0.88}$  \\
 8  & $\mathbf{0.878}$  & $\mathbf{+1.85}$  & $0.245$           & $-1.40$           & $0.847$           & $+1.49$           & $0.194$           & $-1.38$           & $0.776$           & $+1.32$           & $\mathbf{0.255}$  & $\mathbf{-1.34}$  \\
 9  & $\mathbf{0.884}$  & $\mathbf{+1.91}$  & $0.263$           & $\mathbf{-0.95}$  & $0.789$           & $+1.31$           & $0.232$           & $-1.31$           & $0.726$           & $+0.96$           & $\mathbf{0.284}$  & $-1.23$           \\
10  & $\mathbf{0.848}$  & $+1.89$           & $0.269$           & $-1.32$           & $0.823$           & $\mathbf{+2.01}$  & $\mathbf{0.282}$  & $\mathbf{-1.28}$  & $0.788$           & $+1.69$           & $0.256$           & $-1.33$           \\
\hline
\end{tabular}
\end{table}

\section{Curling Physics Simulator}
\label{app:sim}
The transition $s'=\mathrm{Sim}(s,c,a)$ is a JAX physics engine integrated by
semi-implicit Euler (macro step $dt=0.02$\,s, two substeps, up to $1{,}500$
steps, stopping when all stone speeds fall below $v_{\text{stop}}=0.03$\,m/s).
The thrown stone is released at $[-d_{\text{hog-tee}},\,y_0]$ with
$d_{\text{hog-tee}}=28.35$\,m (full sheet) and initial velocity
$v\,[\cos\theta,\sin\theta]$ and spin $\omega$. Dynamics combine:
(i)~linear ice friction with deceleration $a_{\text{linear}}=0.11\,\text{m/s}^2$
opposing motion;
(ii)~a spin-induced curl lateral acceleration
$a_{\text{curl}} = k_{\text{curl}}\,\omega\,\hat v^\perp\,s_{\text{eff}}$
with $k_{\text{curl}}=0.10$ and a speed-saturated factor
$s_{\text{eff}}=c_{\text{cap}}\tanh(v/c_{\text{cap}})$,
$c_{\text{cap}}=2.5\,\text{m/s}$;
(iii)~exponential spin decay at rate $\gamma_{\text{spin}}=0.15\,\text{s}^{-1}$,
with additional damping $c_{\text{spin-contact}}=80$ during contact.
Stone-stone interaction uses smooth penalty contact forces with stiffness
$k_{\text{penalty}}=2.5\!\times\!10^{4}$\,N/m, damping
$c_{\text{damp}}=220$\,kg/s, and a force clamp of $f_{\max}=2\!\times\!10^{4}$\,N
to prevent stiff-spring blow-ups; stone mass is $m=19.96$\,kg and stone
radius is $r=0.145$\,m. Initial overlaps from data quantisation are removed
by an iterative projection. Sheet walls are disabled; the in-play check is
purely geometric. The action ranges (full sheet, realism-fixed) are
\[
v\in[2.20,\,3.01]\,\text{m/s},\quad
\theta\in[-0.1038,\,0.1038]\,\text{rad},\quad
\omega\in[-7.0,\,7.0]\,\text{rad/s},\quad
y_0\in[-0.23,\,0.23]\,\text{m},
\]
and serve as a hard \emph{search-space limit} during target generation:
KR-UCT widening and one-ply sampling never propose outside these bounds, and
the inverse solver also samples within them. The board has 12 stones (six
per team); two are pre-placed and ten are thrown per end; the new stone
fills the thrower's next free slot. The simulator is exposed as a batched
JAX function and is the only transition called during search, target
generation, and evaluation.

\section{Inverse-Action Solver}
\label{app:inverse}
Given a pre-state and an observed post-state, we recover the release
parameters by minimising a position-matching loss with the cross-entropy
method (CEM) \cite{deboer2005tutorial}. Candidate actions in normalised
$[0,1]^4$ are sampled from a Gaussian, simulated forward, and scored by a
team-ownership-aware, permutation-invariant assignment loss: matched live
stones contribute squared position error, while unmatched targets and
spurious extra in-play stones incur capped penalties ($\le 4\,\text{m}^2$
each), so the loss correctly handles stones knocked out of play. CEM
iterates (population $96$, $30$ generations, elite fraction $0.2$,
EMA-smoothed mean/variance with a variance floor, early-stop at residual
$\le 0.1\,\text{m}^2$) in a coarse-then-refine schedule. We retain solutions
that converge (\texttt{solver\_ok}) with a small refined residual to form
the ``realistic'' inverse dataset. Release parameters are $z$-standardised
by the dataset mean/std $(\mu_a,\sigma_a)$, stored with the checkpoint and
used to map policy outputs back to physical actions.

\section{Architecture and Feature Details}
\label{app:arch}
We use the curling-aware graph-feature configuration that the warm-start
encoder was trained with; the same configuration is used by our standalone
value baseline, so the two are directly comparable. The exact feature set
is pinned in the implementation.

\textbf{Nodes (22).} 12 stone-slot nodes (alive iff the slot holds an
in-play stone), plus 9 \emph{landmark} nodes: the button centre at raw
$(750, 800)$, three release-area reference points across the nearer hog
line (left, centre, right) at $\approx 6.4$\,m forward of the button along
the centre line and at lateral offsets $\{-1.22, 0, +1.22\}$\,m, and five
behind-house takeout-boundary points placed $\approx 2.4$\,m behind the
button (just past the back of the house) at lateral offsets $\{-1.83,
-0.91, 0, +0.91, +1.83\}$\,m. The 22nd node is a \emph{global token} that
is prepended to the node tensor at encode time and whose features are the
condition $c$ projected to the hidden dimension; the latent $h$ is read
out as the output embedding of this token. \textbf{Node features} are
6-dimensional: $[x, y, \text{team\_indicator}, \text{is\_live},
\text{is\_landmark}, \text{spare}]$, where \texttt{team\_indicator} is
$0/1$ for stone slots and $0$ for landmarks; the trailing scalar is
unused.

\textbf{Edges (9-d features).} 4 base pairwise scalars
$[\Delta x, \Delta y, \text{distance}, \text{same\_team\_indicator}]$
plus 5 curling-specific scalars:
\textsc{scorability} (unoccluded angular span of the target seen from the
button, computed from board geometry at the current step and broadcast as a
node feature into each incident edge);
\textsc{reach} (curl-arc reachability on release-source edges, pairwise
angular reachability otherwise);
\textsc{score-out} (outgoing-direction compatibility with drawing to the
button);
\textsc{takeout-out} (outgoing-direction compatibility with ejecting an
opponent through the behind-house corridors);
and the composite $\textsc{reach}\cdot\max(\textsc{score-out},
\textsc{takeout-out})$. Edges are not pruned in the current configuration.
The composite scalars are deterministic functions of the live board state
rather than learned, and they recompute every step. They are projected into
a per-head bias and a value modulation at each transformer layer, so the
attention is \emph{edge-biased} rather than purely token-based.

\textbf{Transformer.} 4 layers, hidden width $256$, $4$ attention heads,
dropout $0.1$; the readout is the output embedding of the global token.
Total $\approx 7.3$M parameters.

\textbf{Heads.} The policy head emits a $K=16$ full-covariance Gaussian
mixture over standardised actions, the value head emits a Gaussian
$(\mu_V,\log\sigma_V^2)$, the dynamics head is a residual MLP
$G(h,a_{\text{box}})$, and the auxiliary two-step-return head $\hat r_2$
is a scalar MLP. A BYOL-style target encoder $E_{\bar\theta}$ ($\tau=0.99$
moving-average tracker, no gradient) supplies the consistency target.

\textbf{Warm start.} The trunk loads the matched modules
(\texttt{global\_token}, \texttt{node\_proj}, \texttt{cond\_proj}, and the
transformer layers) from the human-prior policy checkpoint. The policy head
is loaded by default (and then refined by search distillation). The value
head loads the mean and log-variance projection submodules from a
separately trained standalone value model that uses the same graph
configuration. The BYOL target encoder is synced to the trunk before
training begins.

\textbf{Augmentation.} At training time we apply, independently with
probability $0.5$ each, a horizontal flip (mirror stone $x$ across the
centerline; negate the angle, spin, and lateral-offset components of all
action targets) and a team-slot swap (exchange slots $1$--$6$ with
$7$--$12$ and flip $\text{stone\_block}$). These keep value and outcome
targets invariant while doubling the effective data.

\section{Training and Compute}
\label{app:train}
Models are implemented in PyTorch and optimised with Adam, gradient
clipping at max-norm $1.0$, and a per-rank batch size of $256$ (the
curl-arc edge features allocate $O(\text{batch}\cdot\text{stones}^2)$
memory; larger batches exceed device memory). Warm-started parameters
(trunk, policy, value) use a small learning rate
($\approx 1.5\times10^{-5}$); freshly initialised heads use
$3\times10^{-4}$. Mixed precision is disabled (the policy Cholesky
exponentiation is fp32-only under autocast). Training runs data-parallel
across four NVIDIA A10G GPUs with the \texttt{gloo} backend.

Per-horizon collection runs as a separate process on the JAX simulator
(GPU). The one-ply robust selector uses $M=96$ candidate actions per root
with $K=8$ noisy executions averaged; the KR-UCT tree variant uses
$n_{\mathrm{sims}}\approx 96$--$120$ simulations per decision with
progressive-widening constants and a $z$-normalised kernel bandwidth (exact
values in the released config). At $h=10$ the canonical pre-placed
openings are enumerated and the search budget is stratified across the
$(\text{mode},\,\text{guard\_slot})$ groups so each opening receives equal
coverage.

Consolidation reads the union of all per-horizon shards (uniform mix
across horizons within each batch), warm-starts from the per-stage
champion, resets the optimiser, and trains for $\sim 20$ epochs.

\section{Execution-Noise Model}
\label{app:noise}
Execution noise is a diagonal Student-$t$ ($\nu=5$, heavier-tailed than
Gaussian) on the action, with the $v$ and $\theta$ scales taken from the
maximum-likelihood fit of Yee et al.~\cite{yee2016monte} to the World
Curling Federation's \emph{Draw Shot Challenge} dataset, the standard
real-world calibration source for curling execution error. The speed scale
is $\approx 0.0095$\,m/s; the angle scale is speed-dependent, roughly
$0.0012$--$0.0037$\,rad (larger at lower speeds). For \emph{spin} and
\emph{lateral offset} no comparable public calibration data is available,
so we use small empirical constants. We expect a calibrated spin/offset
noise model to be especially relevant for sweeping-influenced trajectories
and near-the-line release offsets; we list this as future work.

The same noise model is used both to perturb realized transitions during
search and evaluation and to average candidate values during target
generation, so the learned policy is optimised for robustness to it.

\section{Free-Guard-Zone Legality}
\label{app:fgz}
With horizon $h$ = throws remaining including the current throw and
$\text{shots\_in\_end}=10$, the thrown-stone index is $k=11-h$. The
no-takeout rule is active for $k\in\{1,2,3\}$ (i.e.\ $h\ge 8$). A post-state
is illegal iff some opponent rock that was in play before the throw is out
of play after it; movement that keeps it in play and removal of one's own
rock are legal. The protected set is all in-play opponent rocks, including
the pre-placed one. Illegal posts are replaced by the pre-shot board
(throw forfeited, displaced rocks restored). The operator is applied to
every simulated candidate post-state in collection, search, and
evaluation.

\bibliographystyle{plain}
\bibliography{references}

\end{document}
